% ==================== 摘要 ====================
{\bfseries 摘要}\quad
本文以 8 组经 200--900 $^\circ$C 热处理的 TiO$_2$ 样品为对象，说明如何用 \texttt{make\_\allowbreak origin\_\allowbreak ready.py} 把仪器导出的 txt 整理成 Origin 可导入的 3 行表头格式，按温度绘制多曲线 XRD 图；并在 Jade 中按化学式检索锐钛矿 71-1167 和金红石 75-1756，导出 $2\theta$ 与相对强度后，在 Origin 中用 Stackplot 画成 PDF 卡片并作为新层叠加到主图。结合本次实验，总结样品制备、开机/关机、校读、扫描和数据整理的操作要点与注意事项，列出低角背底高、峰位偏移、两相峰重叠、Origin 导入错列、Jade 卡片选错等问题及解决办法。200--600 $^\circ$C 以锐钛矿为主，700 $^\circ$C 出现金红石 (110)，800 $^\circ$C 两相共存，900 $^\circ$C 转为金红石。

\vspace{4pt}

\par\noindent{\heiti 关键词：}XRD；TiO$_2$；PDF 卡片；多曲线绘图；Origin；Jade
